% ============================================================================
% Fichier  : partie2/P02_C07_risques.tex
% Titre    : Analyser et Traiter les Risques de Sécurité
% Auteur   : MINKA MI NGUIDJOÏ Thierry Emmanuel
% Version  : 1.0 -- novembre 2026
% ============================================================================
% Base : brouillon « Livre II » (analyse des risques). Cohérence avec INF3226,
% Chapitres 4 et 5 : échelles 1 à 4, score intolérable à partir de 9,
% scénario EBIOS du prestataire de Garoua. Sources : ISO 31000:2018 ;
% ISO/IEC 27005:2022 ; ANSSI, EBIOS Risk Manager (2018) ; Cox (2008).
% ============================================================================

\chapter{Analyser et Traiter les Risques de Sécurité}
\label{chap:risques}
\chaptermark{Risques de sécurité}

\epigraph{%
    « On ne peut pas tout protéger également~:\\
    le risque dit où mettre l'effort. »%
}{Principe du cours}

\niveauChapitre{Fondamental}{Chapitres~\ref{chap:proprietes},
\ref{chap:menaces} et~\ref{chap:gouvernance}.}

\begin{boxobjectifs}
\begin{itemize}
    \item Décrire le processus de gestion des risques de l'ISO~31000 et
          de l'ISO/IEC~27005.
    \item Distinguer l'approche par actifs et l'approche par scénarios.
    \item Conduire les cinq ateliers d'EBIOS Risk Manager sur un cas.
    \item Coter un risque et le placer dans une matrice.
    \item Choisir une option de traitement et rédiger un plan.
    \item Connaître les limites des matrices de risque.
\end{itemize}
\end{boxobjectifs}

\bigskip

% ============================================================================
\section*{Question directrice : le risque de Garoua avait-il été apprécié~?}
% ============================================================================

Non. Aucun document de LiZbiZ ne mentionnait l'accès distant du
prestataire, ni la possibilité de perdre la facturation plusieurs jours.
Ce chapitre montre comment une appréciation des risques, conduite six mois
plus tôt, aurait fait apparaître le scénario de mars et désigné les
mesures qui l'auraient arrêté.

\medskip
\textbf{Plan du chapitre.}
\begin{enumerate}
    \item Le processus de gestion des risques (section~\ref{sec:c07-processus}).
    \item Actifs ou scénarios (section~\ref{sec:c07-approches}).
    \item EBIOS Risk Manager appliqué à LiZbiZ (section~\ref{sec:c07-ebios}).
    \item Coter et cartographier (section~\ref{sec:c07-coter}).
    \item Traiter les risques (section~\ref{sec:c07-traiter}).
    \item Les limites des matrices (section~\ref{sec:c07-limites}).
\end{enumerate}

\separateur

% ============================================================================
\section{Le Processus de Gestion des Risques}
\label{sec:c07-processus}
% ============================================================================

L'ISO~31000 décrit un processus valable pour tout risque~\cite{iso31000_2018}~;
l'ISO/IEC~27005 l'applique à la sécurité de l'information~\cite{iso27005_2022}.

\begin{tableaucours}{L{3.6cm}L{8.6cm}}{%
    \tableauHeader{Étape} &
    \tableauHeader{Contenu}
}
Établir le contexte & Périmètre, objectifs, critères (échelles,
appétence) \\
Identifier & Les sources de risque, les événements, leurs causes et leurs
conséquences \\
Analyser & Estimer la vraisemblance et l'impact \\
Évaluer & Comparer au critère d'acceptation, classer \\
Traiter & Choisir et mettre en œuvre des options \\
Surveiller et revoir & Suivre l'évolution des risques et des mesures \\
Communiquer et enregistrer & À chaque étape, avec les parties
concernées \\
\end{tableaucours}
\vspace{-0.8em}
\begin{center}{\small\itshape Tableau~7.1. Le processus de l'ISO~31000}\end{center}

% ============================================================================
\section{Actifs ou Scénarios}
\label{sec:c07-approches}
% ============================================================================

L'ISO/IEC~27005, dans son édition de 2022, admet deux approches.
L'\termecle{approche par actifs} part de chaque actif, de ses menaces et
de ses vulnérabilités~: elle est exhaustive mais produit des listes
longues. L'\termecle{approche par scénarios} part des événements redoutés
et des chemins d'attaque~: elle est plus proche de la réalité des
attaques. EBIOS Risk Manager, la méthode de l'ANSSI, est une approche par
scénarios~\cite{anssi2018ebios}.

% ============================================================================
\section{EBIOS Risk Manager appliqué à LiZbiZ}
\label{sec:c07-ebios}
% ============================================================================

\begin{tableaucours}{L{2.4cm}L{4.4cm}L{5.4cm}}{%
    \tableauHeader{Atelier} &
    \tableauHeader{Question} &
    \tableauHeader{Résultat chez LiZbiZ}
}
1. Cadrage et socle & Que protège-t-on, et quelles règles applique-t-on
déjà~? & Valeurs métier~: facturation, suivi des expéditions, paiements~;
événement redouté ER1~: facturation arrêtée plus de 3~jours \\
2. Sources de risque & Qui pourrait attaquer, et pour quoi~? & Groupe
criminel (rançon), complice interne (fraude), ancien salarié (revente de
données) \\
3. Scénarios stratégiques & Par quel chemin dans l'écosystème~? & Par le
prestataire de Garoua, qui a un accès distant \\
4. Scénarios opérationnels & Par quelles étapes techniques~? &
Hameçonnage, accès distant sans second facteur, serveur non corrigé,
sauvegardes en ligne, chiffrement \\
5. Traitement & Quelles mesures, quel risque résiduel~? & Double
authentification, correctifs sous 30~jours, sauvegardes hors ligne,
clauses au contrat \\
\end{tableaucours}
\vspace{-0.8em}
\begin{center}{\small\itshape Tableau~7.2. Les cinq ateliers d'EBIOS RM chez
LiZbiZ}\end{center}

\begin{boxscenario}{Opération GAROUA. Le scénario était écrit d'avance}
\textbf{Constatation.} Le scénario opérationnel de l'atelier~4, que le
cours d'audit (UE INF3226) construit pour LiZbiZ, décrit presque mot pour
mot l'attaque de mars~: hameçonnage d'un technicien du prestataire, accès
distant sans double authentification, élévation de privilèges sur un
serveur non corrigé, destruction des sauvegardes en ligne, chiffrement.

\medskip
\textbf{Analyse} (confiance élevée). L'attaque n'était pas imprévisible~;
elle n'avait pas été prévue. Une appréciation des risques en septembre,
au moment de confier les centres à SAHEL INFOGÉRANCE, aurait suffi.
\end{boxscenario}

% ============================================================================
\section{Coter et Cartographier}
\label{sec:c07-coter}
% ============================================================================

On cote chaque risque en vraisemblance (V) et en impact (I), de 1 à 4,
selon des échelles fixées par la direction. Le score est \(V \times I\),
de 1 à~16.

\begin{tableaucours}{cL{5.2cm}L{5.8cm}}{%
    \tableauHeader{Niveau} &
    \tableauHeader{Vraisemblance (sur un an)} &
    \tableauHeader{Impact}
}
1 & Rare & Mineur~: perte de moins de 5~millions FCFA, arrêt de moins de
4~h \\
2 & Possible & Modéré~: de 5 à 50~millions, arrêt de moins de 24~h \\
3 & Probable & Majeur~: de 50 à 500~millions, arrêt de 1 à 3~jours \\
4 & Quasi certain & Critique~: plus de 500~millions, arrêt de plus de
3~jours, sanction \\
\end{tableaucours}
\vspace{-0.8em}
\begin{center}{\small\itshape Tableau~7.3. Échelles de cotation (convention du cours)}\end{center}

\begin{tableaucours}{L{1.0cm}L{5.6cm}cccL{2.0cm}}{%
    \tableauHeader{Réf.} &
    \tableauHeader{Risque} &
    \tableauHeader{V} &
    \tableauHeader{I} &
    \tableauHeader{Score} &
    \tableauHeader{Zone}
}
S1 & Rançongiciel par l'accès distant du prestataire & 3 & 4 & 12 &
Intolérable \\
S2 & Vol des données clients & 3 & 3 & 9 & Intolérable \\
S3 & Fraude au virement par hameçonnage ciblé & 2 & 3 & 6 & À traiter \\
S4 & Panne du centre de Garoua & 2 & 3 & 6 & À traiter \\
S5 & Défiguration du portail de suivi & 2 & 2 & 4 & Acceptable \\
\end{tableaucours}
\vspace{-0.8em}
\begin{center}{\small\itshape Tableau~7.4. Registre des risques de sécurité de
LiZbiZ (septembre, avant traitement)}\end{center}

Le conseil a fixé l'appétence~: un score de 9 ou plus est intolérable et
traité sans délai~; de 5 à 8, il est traité dans l'année~; de 1 à 4, il
est accepté sous surveillance.

% ============================================================================
\section{Traiter les Risques}
\label{sec:c07-traiter}
% ============================================================================

\begin{tableaucours}{L{2.4cm}L{5.4cm}L{4.4cm}}{%
    \tableauHeader{Option} &
    \tableauHeader{Principe} &
    \tableauHeader{Exemple chez LiZbiZ}
}
Réduire & Mettre en place des mesures & Double authentification,
sauvegardes hors ligne \\
Transférer & Partager les conséquences financières & Assurance
cyber~; le risque juridique reste \\
Éviter & Renoncer à l'activité à risque & Supprimer l'accès distant
permanent du prestataire \\
Accepter & Décider en connaissance de cause, par écrit & S5, sous
surveillance \\
\end{tableaucours}
\vspace{-0.8em}
\begin{center}{\small\itshape Tableau~7.5. Les options de traitement}\end{center}

\begin{boxobjet}{M-07 · Sauvegardes hors ligne testées}
\textbf{Objectif~:} pouvoir reprendre l'activité après un chiffrement
malveillant. \textbf{Risque traité~:} S1. \textbf{Mise en œuvre~:} une
copie quotidienne sur un support déconnecté ou immuable, hors du domaine,
conservée 30~jours. \textbf{Vérification~:} restauration complète d'un
serveur chaque trimestre, chronométrée et consignée.
\end{boxobjet}

\begin{boxalert}
\textbf{Violation typique.} Abaisser la vraisemblance d'un risque dès
qu'une mesure est \emph{décidée}.

\smallskip
\textbf{Règle~:} le risque résiduel ne tient compte que des mesures en
place \emph{et dont l'efficacité est vérifiée}. Une sauvegarde jamais
restaurée ne réduit aucun risque.
\end{boxalert}

% ============================================================================
\section{Les Limites des Matrices}
\label{sec:c07-limites}
% ============================================================================

L.~Cox a montré que les matrices de risque peuvent classer des risques
dans le mauvais ordre~: deux risques de même score peuvent avoir des
pertes attendues très différentes, et le produit de deux rangs n'est pas
une grandeur~\cite{cox2008matrices}. Les matrices servent à trier et à
dialoguer~; pour arbitrer un budget, on complète par une estimation
chiffrée, en intervalles (méthode FAIR, cours INF3226).

\begin{boxmetier}{RSSI}
\textbf{Tenir le registre.} Le RSSI tient le registre des risques~: pour
chacun, son propriétaire, sa cote avant et après traitement, les mesures,
leur état de vérification, et la décision d'acceptation signée. Il le
revoit à chaque changement majeur (nouveau prestataire, nouvelle
application) et après chaque incident.
\end{boxmetier}

\begin{boxmetier}{Auditeur de sécurité}
\textbf{Tester le registre.} L'auditeur vérifie que chaque mesure qui
réduit un risque a été testée~; sinon, il recalcule le risque résiduel
sans elle.
\end{boxmetier}

\separateur

\begin{boxresume}
\textbf{Ce qu'il faut retenir de ce chapitre~:}
\begin{itemize}
    \item \textbf{Le processus} (ISO~31000, ISO/IEC~27005)~: contexte,
          identification, analyse, évaluation, traitement, surveillance.
    \item \textbf{L'approche par scénarios} (EBIOS RM, cinq ateliers) colle
          aux attaques réelles~; chez LiZbiZ, elle décrivait l'attaque de
          mars.
    \item \textbf{La cotation} \(V \times I\) classe les risques au regard
          de l'appétence fixée par la direction.
    \item \textbf{Quatre options}~: réduire, transférer, éviter, accepter~;
          le résiduel ne compte que les mesures vérifiées.
    \item \textbf{Les matrices} trient~; elles ne suffisent pas pour
          arbitrer un budget.
\end{itemize}
\end{boxresume}

\begin{boxinfo}
\textbf{Points de contrôle de l'élève.} Je sais~: décrire le processus de
l'ISO~31000~; distinguer approche par actifs et par scénarios~; dérouler
les cinq ateliers d'EBIOS RM~; coter un risque et le classer selon
l'appétence~; choisir une option de traitement~; citer deux limites des
matrices.
\end{boxinfo}

\bigskip

\begin{boxexo}[title={\ding{45}\quad Exercice~7.1. Le processus \niveaubase}]
Rangez dans l'ordre du processus~: comparer au critère d'acceptation~;
fixer les échelles~; installer la double authentification~; estimer la
vraisemblance~; recenser les événements redoutés.
\end{boxexo}

\begin{boxexo}[title={\ding{45}\quad Exercice~7.2. Coter \niveaubase}]
Calculez le score et la zone de~: (a) V~=~4, I~=~2~; (b) V~=~1, I~=~4~;
(c) V~=~3, I~=~3.
\end{boxexo}

\begin{boxexo}[title={\ding{45}\quad Exercice~7.3. Quelle option~? \niveaubase}]
Associez une option de traitement à~: souscrire une assurance~; fermer le
portail de suivi au public~; chiffrer les sauvegardes~; garder un
serveur obsolète isolé en attendant son remplacement, sur décision écrite.
\end{boxexo}

\begin{boxexo}[title={\ding{45}\quad Exercice~7.4. Atelier 2 \niveaumoyen}]
Conduisez l'atelier~2 d'EBIOS RM pour la BANQUE SANAGA~: au moins quatre
sources de risque, leurs objectifs, et lesquelles retenir.
\end{boxexo}

\begin{boxexo}[title={\ding{45}\quad Exercice~7.5. Scénario opérationnel \niveaumoyen}]
Écrivez le scénario opérationnel du risque S3 (fraude au virement par
hameçonnage ciblé) en cinq étapes, et une mesure par étape.
\end{boxexo}

\begin{boxexo}[title={\ding{45}\quad Exercice~7.6. Risque résiduel \niveaumoyen}]
Après traitement, la double authentification est en place et testée~;
les sauvegardes hors ligne sont achetées mais jamais restaurées. Recotez
S1 et justifiez.
\end{boxexo}

\begin{boxexo}[title={\ding{45}\quad Exercice~7.7. Transférer \niveaumoyen}]
Une assurance cyber couvre 300~millions FCFA de pertes. Quels éléments du
risque S1 reste-t-il à LiZbiZ~?
\end{boxexo}

\begin{boxexo}[title={\ding{45}\quad Exercice~7.8. Les limites \niveaumoyen}]
Deux risques ont le score 8~: V~=~2, I~=~4 et V~=~4, I~=~2. Donnez un
exemple chiffré où leurs pertes annuelles attendues diffèrent d'un
facteur 10, et tirez-en la leçon.
\end{boxexo}

\begin{boxexo}[title={\ding{45}\quad Exercice~7.9. Le registre \niveauexpert}]
Construisez le registre des risques de sécurité de l'ANPP (cinq risques),
avec cotation, zone, option, mesures et propriétaire.
\end{boxexo}

\begin{boxexo}[title={\ding{45}\quad Exercice~7.10. Septembre \niveauexpert}]
Vous êtes le RSSI de LiZbiZ en septembre, au moment du contrat avec SAHEL
INFOGÉRANCE. Rédigez la note d'une page au directeur général qui expose
le risque S1, trois mesures et leur coût approximatif, et demande une
décision.
\end{boxexo}

\bigskip

\begin{boxinfo}
\textbf{Pour aller plus loin.}
\begin{itemize}
    \item ISO~31000 et ISO/IEC~27005~\cite{iso31000_2018,iso27005_2022}.
    \item ANSSI, \emph{EBIOS Risk Manager}~\cite{anssi2018ebios}.
    \item L.~A.~Cox, sur les matrices de risque~\cite{cox2008matrices}.
\end{itemize}
\end{boxinfo}

\begin{boxinfo}
\textbf{Transition.} Une appréciation des risques suppose de connaître
ses actifs. Le Chapitre~\ref{chap:actifs} montre comment les
inventorier, les classer et les protéger tout au long de leur vie.
\end{boxinfo}
